We present the first analytically tractable instance of the spectral admissibility
programme proposed in the Cosmochrony framework.
On the undirected Cayley graph of the quaternion group $Q_8$, with generators
$S=\{\pm i,\pm j,\pm k\}$, a local bounded–flux constraint
$|\partial_t\chi_v|\le c_\chi$ induces a hierarchy of maximal admissible
mode amplitudes across irreducible spectral sectors.

In the linearised regime, the admissibility bound takes the universal form
\[
  A_n^{\max}=\frac{c_\chi}{\sqrt{\lambda_n}},
\]
where $\lambda_n$ are the Laplacian eigenvalues of the graph.
This produces a spectral hierarchy in which the non-abelian sector
admits a larger maximal effective amplitude than the nontrivial abelian sectors
by a factor $\sqrt{4/3}$.

We then perform a weakly nonlinear modal reduction of the discrete
Dirac–Born–Infeld dynamics and show that the first nonlinear corrections
preserve this hierarchy.
Solving the one-mode DBI dynamics exactly on the $Q_8$ graph, we further show
that the same envelope $A_n^{\max}=c_\chi/\sqrt{\lambda_n}$ coincides with the
exact saturation boundary, demonstrating structural stability beyond the
perturbative regime.

The analysis is extended to the binary icosahedral group $2I$, the largest finite
subgroup of $\mathrm{SU}(2)$, where the same admissibility envelope generates
a four-level hierarchy across spin sectors with contrast $\sqrt{14/9}$.
Together these results provide explicit realizations of spectral admissibility
along the chain
\[
  Q_8 \subset 2I \subset \mathrm{SU}(2).
\]

More generally, the results show that bounded substrate capacity can act as a
purely spectral filter favouring low-lying non-abelian sectors of relational
dynamics.

\medskip
\noindent\textbf{Keywords:} Spectral graph theory, Cayley graphs, Representation theory,
Laplacian spectrum, Dirac–Born–Infeld dynamics, Non-abelian symmetry, Bounded-flux dynamics
